\textbf{Ramal de potencia vigente INC-42.} Se sustituye el límite de cinco metros por montaje adyacente FRPS/VEP con tramo máximo de un metro, cable seleccionado en INC-42 y T-11/48L. El presupuesto y caída de cinco metros que siguen son antecedentes, sin representar otro ramal instalado. Se conservan fuentes, detectores, EOLR y monitores; autonomía útil y subtensión en bornes continúan pendientes.

\section{Supervisión elegida del ramal de aspiración ordinaria (INC-39)}
\label{sec:sd4-inc39}
Se mantienen seis VEP, cuatro FRPS, ocho baterías de 33 Ah, cuatro gabinetes de esas baterías y los dos pares de 75 Ah del panel/SPS. Para cada VEP ordinario F01/F02/F03/F06 se selecciona un relé System Sensor EOLR-1, referencia Fike 02-4981, instalado en caja ordinaria seca junto a los bornes de alimentación del detector. Son cuatro activos y uno de reserva fría. El manual completo FRPS 06-583 rev.~3 exige ese relé UL de final de línea cuando una salida NAC se configura como AUX; la supervisión de la fuente no basta para acreditar el ramal. P7 AUX también exige relé si el equipo no supervisa su conexión. Por ello no se atribuye supervisión continua del cable por conservar el nombre NAC \parencite[secs.~2.5.6.2 y 2.5.8, pp.~19/21]{lote39-frpsmanual}.

Se fija salida 1 de cada FRPS en modo 24 V Power Supply, SW2-4/5/6 OFF; salida 2 del mismo par también en modo fuente continua, SW2-1/2/3 OFF, sin carga conectada. La combinación fuente/fuente está admitida; las salidas 3/4 se mantienen apagadas y P7 no alimenta el VEP. El borne de salida 1 recibe el ramal y el relé al extremo. Es alimentación continua no rearmable, sin orden NAC de sincronización, retenedor de puerta ni corte a treinta segundos por pérdida AC. El rango AUX real de esa salida es 18,6--28,5 V; configurarla como fuente no transforma su mínimo en 24 V regulados. El límite total externo sigue siendo 1 A en reposo, aunque una salida admita 3 A \parencite[secs.~1.3, 3.1 y 3.1.2, pp.~9--10/22--24]{lote39-frpsmanual}.

El EOLR-1 tiene entrada 9--40 V DC, consumo máximo 20 mA, contacto SPST normalmente abierto, 30 V DC/3 A resistivo, 23 por 42 por 31 mm, ambiente -30/+60 grados C y humedad sin condensación. Se conectan rojo/negro a los mismos bornes finales que alimentan VEP, después del cable; los dos violetas proporcionan contacto libre de potencial hacia un monitor 55-046 separado, sin puentear el negativo FRPS con SLC. El contacto cierra mientras tiene alimentación. Se configura entrada de supervisión de circuito sano con EOL 39 k$\Omega$ 10-2625 en el extremo y esquema OEM de detección de cortocircuito 10-2530 cuando corresponda; apertura de ramal/relé genera avería enclavada, no fuego ni autorización de descarga. La apertura física del circuito de monitorización conserva su avería propia. El esquema y valores se reciben contra 06-356/06-457, sin cambiar los resistores para acomodar corrientes \parencite{lote39-eolr1}\parencite[sec.~3.9.2 y ap.~C]{lote33-cheetah11}.

La señal nueva identifica pérdida de alimentación al extremo; no sustituye Fault VEP ni P3 Trouble FRPS. Se conservan ambos estados separados, incluido fallo de fuente durante alarma, porque el contacto interno P6 no garantiza informe de todos los fallos cuando está activo. Se seleccionan cuatro 55-046 y cuatro EOL 10-2625 adicionales a los 52 anteriores: 56 monitores y 128 direcciones activas, conservando 54 relés y sus reservas. Se asignan dos nuevas entradas a SLC3 y dos a SLC4: reparto 40/40/24/24 direcciones, corriente de alarma 88/88/52/52 mA, dentro de 100 mA por lazo. El incremento se justifica para distinguir fuente, ramal y detector, sin usar los dos sensores pendientes de bancos como entradas disponibles.

El relé EOLR-1 puede seguir activado muy por debajo de los 18 V mínimos del VEP; su contacto no es monitor de subtensión a 18 V ni prueba de detector operativo. La aceptación comprueba pérdida de positivo y negativo al extremo, contacto pegado, cortocircuito, conexión floja, P3 FRPS, alimentación sólo por batería y estado Fault VEP. Se registra como evidencia prevista un informe de supervisión y tensión del ramal; ninguna de estas pruebas se presenta como ejecutada. La aptitud de confirmadores independientes de bancos, clasificación y ASPIRE conserva el expediente de INC-36.

\section{Presupuesto actualizado y condición de tensión (EXT-39)}
\label{sec:sd4-ext39}
A 24 V, cada FRPS con VEP y EOLR-1 demanda como subtotal $0,57+0,020+0,035=0,625$ A en reposo y $0,59+0,020+0,141=0,751$ A en alarma. El cribado de 24 h más cinco minutos con factor 1,2 resulta $1,2[24(0,625)+(5/60)(0,751)]=18,0751$ Ah por fuente; se conservan los pares de 33 Ah. Las cuatro suman 60,00/72,096 W DC a 24 V, sin atribuirlos a entrada AC ni factor de potencia. El EOLR-1 agrega 0,48 W por fuente en ese punto. Los máximos VEP 0,57/0,59 A están publicados a 24 V; no se extrapolan como máximos de corriente al extremo inferior del rango AUX \parencite{lote39-eolr1}\parencite[p.~2]{lote30-vep36106}.

Los cuatro monitores adicionales agregan 1,940 mA en reposo y 8 mA en alarma SLC. El subtotal del controlador pasa de 1,230746/5,768 A a 1,232686/5,776 A, y su cribado a $1,2[24(1,232686)+(5/60)(5,776)]=36,078957$ Ah. SPS conserva 0,7954/3,45 A y 23,2525 Ah. No se trasladan consumos de bobinas EOLR a Cheetah: se alimentan de FRPS. Estos subtotales no cierran autonomía útil, consumo VEP por tensión, tiempo de recarga o capacidad de baterías a temperatura/tensión final.

Se fija un límite de instalación propio de cinco metros de ida entre cada FRPS y su VEP ordinario, sin sumar al recorrido otros detectores. El ramal usa cobre multifilar AWG14 y resistencia de cribado recubierta de 10,7 $\Omega$/km a 20 grados C del manual Fike; factor térmico propio a 55 grados C de 1,13755 y reserva de 0,10 V para bornes. La corriente externa admisible se limita en el diseño a 1 A total, incluidos VEP/EOLR y restantes salidas. En ese extremo, $\Delta V\leq2(0,005)(10,7)(1,13755)(1,0)+0,10=0,221718$ V; con 18,6 V al origen, el extremo sería al menos 18,378282 V, margen 0,378282 V sobre el mínimo VEP. A 28,5 V no se supera su máximo 30 V. Son desigualdades de proyecto, sujetas a mínimo real de salida en todos los estados y corriente real VEP a esa tensión; el fusible no demuestra que una carga operativa se mantenga bajo 1 A sin dispararse \parencite[sec.~1.3]{lote39-frpsmanual}\parencite[p.~4-3]{lote33-cheetah11}.

El manual FRPS establece retirada de salidas cuando la batería baja de 18,4 V. Ese corte no prueba mantener 18,6 V al origen durante toda la autonomía; tampoco justifica corriente VEP a batería baja. La memoria final coteja curva útil del par 33 Ah, regulación/corte, consumo a 18--30 V, recarga, cable/conectores y tensión medida en el VEP con alarma y sólo batería. Una colocación que exceda cinco metros requiere recalcular o reubicar fuente en pared ordinaria; no invade rutas ni bancos para cumplir longitud. La caída DFIA de EXT-36 conserva su condición de mínima fuente 22 V no acreditada; no se confunde el rango de FRPS con la fuente de liberación Cheetah.

\section{Identificación de ramales y conservación de compartimentación (PEN-39)}
\label{sec:sd4-pen39}
Se identifican cuatro ramales nuevos de alimentación/supervisión de extremo: INC-F01-PWR/EOL, INC-F02-PWR/EOL, INC-F03-PWR/EOL e INC-F06-PWR/EOL. Cada FRPS, VEP, EOLR y monitor tiene un ID propio; el esquema relaciona los dos conductores de potencia con la pareja independiente del contacto supervisado. El relé queda en caja ordinaria accesible junto al VEP, antes de cualquier tramo final que pudiera quedar sin vigilancia; la conexión a bornes finales se recibe visualmente. No se instala EOLR ordinario en bancos ni se utiliza su intervalo ambiental como aprobación H$_2$.

La instalación conserva treinta aberturas y sesenta SF/Wedge, con lista de empaquetado por ambas caras. El nuevo calibre AWG14 y la longitud máxima de cinco metros son parámetros de cálculo, no diámetro exterior ni cubierta de un cable comercial. Se mantiene pendiente elegir variante de cable con ficha y cotejo del listado C-AJ-8309; la pareja de contacto y la alimentación no se introducen como un nuevo servicio certificado por existir un marco SF. Los ramales van por control y sus pasos ya reservados, separados de energía A/B, sin agregar cuatro marcos por deducción ni compartir los recorridos de la fuente de liberación. Piso/adaptación, enlaces dinámicos marinos, juntas y soporte compuesto conservan sus cálculos y verificación \parencite[pp.~3--4]{lote30-caj8309}.

La supervisión seleccionada desarrolla una brecha concreta de energía de detección; no completa RT-06.04, RT-06.16 ni RT-06.17. Se conservan selección láser, límites de tensión/autonomía y obligaciones de confirmación independiente, certificado íntegro, clasificación, gas residual, hidráulica, concentración, retención y alivio. No se convierte un relé de pérdida de alimentación en barrera intrínsecamente segura ni se habilita descarga/sellado por haber supervisado un ramal.
